\item[\textbf{Exercise 15.3-1}]
Explain why, in the proof of Lemma 15.2, if $\mathit{x.freq} = \mathit{b.freq}$,
then we must have $\mathit{a.freq} = \mathit{b.freq} = \mathit{x.freq} =
\mathit{y.freq}$. 

\Solution

From the premises of Lemma 15.2, we know that $x$ and $y$ are the two characters
with the lowest frequencies, meaning $x.freq \le y.freq \le c.freq$ for any
other character $c$. We are also given two sibling characters $a$ and $b$ at the
maximum depth, where $a.freq \le b.freq$. 

Because $x$ and $y$ have the lowest frequencies overall, it must be true that
$x.freq \le a.freq$ and $y.freq \le b.freq$. 

Given the assumption that $\mathit{x.freq} = \mathit{b.freq}$, we can prove the
equality for $a$:
\begin{alignat*}{2}
                      & \mathit{x.freq} \le \mathit{a.freq} \le \mathit{b.freq} \quad && \text{(From the definitions above)} \\
    \Rightarrow \quad & \mathit{x.freq} \le \mathit{a.freq} \le \mathit{x.freq} \quad && \text{(Substituting $\mathit{b.freq}$ with $\mathit{x.freq}$)} \\
    \Rightarrow \quad & \mathit{a.freq} = \mathit{x.freq} = \mathit{b.freq} &&
\end{alignat*}

Similarly, we can prove the equality for $y$:
\begin{alignat*}{2}
                      & \mathit{x.freq} \le \mathit{y.freq} \le \mathit{b.freq} \quad && \text{(From the definitions above)} \\
    \Rightarrow \quad & \mathit{x.freq} \le \mathit{y.freq} \le \mathit{x.freq} \quad && \text{(Substituting $\mathit{b.freq}$ with $\mathit{x.freq}$)} \\
    \Rightarrow \quad & \mathit{y.freq} = \mathit{x.freq} = \mathit{b.freq} &&
\end{alignat*}

By combining these two proven equalities, we get the final equation:
$$
    \mathit{x.freq} = \mathit{y.freq} = \mathit{b.freq} = \mathit{a.freq}
$$
